% Part: computability
% Chapter: recursive-functions
% Section: primes

\documentclass[../../../include/open-logic-section]{subfiles}

\begin{document}

\olfileid{cmp}{rec}{pri}
\olsection{الأعداد الأولية}

يوفر لنا التكميم المحدود والتصغير المحدود أدوات وافرة لبيان أن الدوال
والعلائق الطبيعية عودية بدائية. انظر، مثلًا، في العلاقة «$x$ يقسم
$y$»، التي تكتب $x\mid y$. تصدق العلاقة $x\mid y$ إذا أمكنت قسمة
$y$ على~$x$ بلا باق، أي إذا كان $y$ مضاعفًا صحيحًا لـ~$x$. (وإذا لم
تصدق، أي إذا كان باقي قسمة $y$ على $x$ أكبر من~$0$، كتبنا
$x\nmid y$.) وبعبارة أخرى، $x\mid y$ إذا وفقط إذا وجد~$z$ بحيث
$x\cdot z=y$. ومن الواضح أن كل $z$ كهذا، إن وجد، يجب أن يكون
$\leq y$. ومن ثم يصدق $x\mid y$ إذا وفقط إذا وجد $z\le y$ بحيث
$x\cdot z=y$. ويمكننا تعريف العلاقة $x\mid y$ بالتكميم الوجودي
المحدود من المساواة والضرب هكذا:
\[
x \mid y \defiff \bexists{z \leq y}{(x \cdot z) = y}.
\]
وهكذا بينا أن $x\mid y$ عودية بدائية.

يكون العدد الطبيعي~$x$ \emph{أوليًا} إذا لم يكن $0$ ولا $1$، ولم
يكن قابلًا للقسمة إلا على $1$ وعلى نفسه. وبعبارة أخرى، تكون الأعداد
الأولية بحيث إذا كان $y\mid x$، فإما $y=1$ وإما~$y=x$. ولاختبار ما
إذا كان $x$ أوليًا، لا يلزمنا إلا فحص $y\mid x$ لكل $y\le x$، لأن
$y\nmid x$ تلقائيًا إذا كان $y>x$. ولذلك يمكن تعريف العلاقة
$\fn{Prime}(x)$، التي تصدق إذا وفقط إذا كان $x$ أوليًا، بـ
\[
\fn{Prime}(x) \defiff x \geq 2 \land \bforall{y \leq x}{(y \mid x \lif y
  = 1 \lor y = x)}
\]
فتكون عودية بدائية.

الأعداد الأولية هي $2$ و$3$ و$5$ و$7$ و$11$، وهكذا. انظر في الدالة
$p(x)$ التي تعيد العدد الأولي ذا الرتبة $x$ في هذه المتتالية؛ أي
$p(0)=2$ و$p(1)=3$ و$p(2)=5$، وهكذا. (وتيسيرًا للكتابة سنكتب غالبًا
$p_x$ بدلًا من $p(x)$؛ فـ$p_0=2$ و$p_1=3$، وهكذا.)

لو كانت لدينا دالة $\fn{nextPrime(x)}$ تعيد أول عدد أولي أكبر من~$x$،
لأمكن تعريف $p$ بسهولة بالعودية البدائية:
\begin{align*}
  p(0) & = 2\\
  p(x+1) & = \fn{nextPrime}(p(x))
\end{align*}
ولما كانت $\fn{nextPrime}(x)$ أصغر $y$ بحيث $y>x$ و$y$ أولي، أمكن
حسابها بسهولة بالبحث غير المحدود. غير أنه يمكن تعريفها أيضًا بالتصغير
المحدود، بفضل نتيجة ترجع إلى إقليدس: يوجد دائمًا عدد أولي بين $x$ و
$\fact{x}+1$.
\[
  \fn{nextPrime(x)} =
  \bmin{y \leq \fact{x}+1}{(y > x \land \fn{Prime}(y))}.
\]
يبين هذا أن $\fn{nextPrime}(x)$، ومن ثم $p(x)$، ليستا قابلتين للحساب
فحسب، بل عوديتان بدائيتان.

(وإن شئت برهانًا سريعًا لمبرهنة إقليدس، فإليك إياه. افترض أن $p_n$
أكبر عدد أولي $\le x$، وانظر في الجداء
$p=p_0\cdot p_1\cdot\dots\cdot p_n$ لجميع الأعداد الأولية~$\le x$.
إما أن يكون $p+1$ أوليًا، وإما أن يوجد عدد أولي بين $x$ و~$p+1$.
لماذا؟ افترض أن $p+1$ ليس أوليًا. عندئذ يوجد عدد أولي $q\mid p+1$
بحيث $q<p+1$. ولا يقسم أي عدد أولي $\le x$ العدد $p+1$. (فبحسب
تعريف~$p$، يقسم كل عدد أولي $p_i\le x$ العدد~$p$ بلا باق. ولذلك تكون
قسمة $p+1$ على كل $p_i\le x$ ذات باق~$1$، ومن ثم
$p_i\nmid p+1$.) وهكذا يكون $q$ عددًا أوليًا أكبر من~$x$ وأصغر من
$p+1$. ولأن $p\le\fact{x}$، يوجد عدد أولي أكبر من~$x$ وأصغر من أو
مساو لـ$\fact{x}+1$.)

\begin{prob}
عرف القسمة الصحيحة $d(x,y)$ باستعمال التصغير المحدود.
\end{prob}

\end{document}
